\documentclass{article}

\usepackage[utf8]{inputenc}
\usepackage{hyperref}
\usepackage{booktabs}
\usepackage{enumitem}
\usepackage{titlesec}
\usepackage{microtype}

\hypersetup{colorlinks=true,linkcolor=black,urlcolor=blue,citecolor=black}

\title{Infosec OS.1:\\
High-Assurance Architecture on Synrc Hypervision}
\author{Namdak Tonpa\\Synrc Research}
\date{August 2026}

\begin{document}

\maketitle

\begin{abstract}
Infosec OS.1 is a high-assurance operating system built on Synrc Hypervision,
a hybrid Type-1 architecture that hosts an unmodified Erlang/OTP BEAM runtime
on the formally verified seL4 microkernel while confining device drivers to a
minimal Alpine Linux guest. The system is constructed entirely statically under
the Microkit framework. This article presents the architecture, derives
technical requirements according to ISO/IEC/IEEE 29148:2018, maps selected
NIST SP 800-53 control families onto seL4 capabilities, and shows how elevating
key POSIX applications (BE, TV, SH, SKYNET) into dedicated Protection Domains
further strengthens isolation, least privilege, and auditability.
\end{abstract}

\tableofcontents
\newpage

%----------------------------------------------------------------------
\section{Introduction}
%----------------------------------------------------------------------

Erlang/OTP and its BEAM virtual machine provide robust concurrency, supervision
trees, and soft real-time behaviour that are attractive for high-availability
and distributed systems. These strengths, however, traditionally rest on a
large general-purpose operating system whose trusted computing base (TCB)
undermines strong security and certification arguments.

Formal isolation has advanced substantially with the seL4 microkernel, which
offers machine-checked proofs of functional correctness, integrity, and
(for selected configurations) information-flow security. Running a full Linux
guest under seL4 merely relocates the TCB problem; pure unikernel approaches
often force substantial rewriting of application runtimes or sacrifice device
support.

Synrc Hypervision resolves this tension by combining three elements:
\begin{itemize}[leftmargin=*]
  \item the formally verified seL4 microkernel and its Microkit framework for
        static, capability-based system construction;
  \item Synrc, a minimal host that supplies only the approximately fifty
        Linux-compatible system calls required by a musl-linked, unmodified
        OTP BEAM;
  \item a carefully confined Alpine Linux virtual machine that supplies real
        device drivers, communicating with the BEAM domain solely through
        capability-mediated shared-memory channels defined by the seL4 Device
        Driver Framework (sDDF).
\end{itemize}

The resulting architecture keeps the TCB under the BEAM extremely small
(seL4 + Synrc) while still permitting practical use of complex hardware.
All isolation boundaries, scheduling parameters, and device ownership are
fixed at image-construction time, yielding a static system amenable to
analysis and aligned with selected families of NIST SP 800-53 controls.

Infosec OS.1 extends this foundation into a coherent security-oriented stack
that also includes a NuttX-based POSIX real-time layer (RT), a drone-swarm
coordination protocol (SKYNET), and a set of zero-dependency Infosec tools
(Binary Editor, Terminal Vision, and a formally modelled POSIX shell).
The present article describes the overall architecture, formalises the
technical requirements, maps the relevant security controls, and demonstrates
the security gains obtained by placing the principal POSIX applications into
their own Protection Domains.

%----------------------------------------------------------------------
\section{Architecture of Infosec OS.1}
%----------------------------------------------------------------------

Infosec OS.1 is organised as a static collection of seL4 Protection Domains
(PDs) whose memory regions, communication channels, interrupt bindings, and
Mixed-Criticality Scheduling (MCS) contexts are declared in a Microkit system
description and materialised at build time. No protection domain or channel
can be created dynamically at runtime.

The principal components are:

\begin{center}
\begin{tabular}{@{}llp{6.5cm}@{}}
\toprule
\textbf{Component} & \textbf{Type} & \textbf{Responsibility} \\
\midrule
seL4               & Microkernel      & Isolation, capabilities, MCS scheduling, IRQ routing \\
Microkit           & Framework        & Static system description and capability layout \\
Synrc              & PD               & Minimal syscall surface and host for BEAM \\
BEAM / OTP         & Runtime          & Erlang/Elixir processes, schedulers, supervision \\
VMM (libvmm)       & PD               & Creates and manages the Alpine Linux guest \\
Alpine Linux       & Guest VM         & Real device drivers (network, block, \ldots) \\
sDDF               & Protocol         & Shared-memory rings and notifications \\
Console PD         & PD               & Exclusive ownership of UART / virtio-console \\
Monitor PD         & PD               & Metrics, health, and audit events \\
Storage / Network  & PD               & Block and Ethernet access (native or mediated) \\
BE, TV, SH, SKYNET & PD               & Isolated POSIX applications (see \S\ref{sec:posix-pds}) \\
\bottomrule
\end{tabular}
\end{center}

All authority is expressed by seL4 capabilities granted at system construction.
There is no ambient authority. Scheduling follows the seL4 MCS model; typical
priority ordering (highest to lowest) places the Console PD first, followed by
Synrc/BEAM, the Monitor, the VMM/Alpine domain, and background domains.
Core affinity and budget/period parameters are supplied via compile-time flags.

\subsection{Technical Requirements (ISO/IEC/IEEE 29148:2018)}
\label{sec:req}

Technical requirements are derived and organised according to the principles of
ISO/IEC/IEEE 29148:2018 (Systems and software engineering --- Life cycle
processes --- Requirements engineering). They are stated as necessary
properties of any conforming implementation of Infosec OS.1.

\begin{enumerate}[leftmargin=*,label=\textbf{TR-\arabic*}]
  \item \textbf{Static construction.}  
        All Protection Domains, memory regions (with explicit read/write/execute
        rights), channels, interrupt bindings, and MCS scheduling contexts shall
        be declared in a Microkit (or equivalent) system description and fixed at
        image-construction time. Dynamic creation of domains or channels is
        prohibited.

  \item \textbf{Minimal syscall surface.}  
        The host for the primary application runtime (Synrc) shall implement only
        the Linux-compatible system calls required by a musl-linked OTP binary
        (memory management, threading and futex family, time, basic file
        descriptors, and an in-memory VFS). The host shall never own real device
        registers.

  \item \textbf{Capability-mediated I/O.}  
        All device access shall occur exclusively through sDDF-style shared-memory
        rings and notifications, or through native seL4 driver PDs. The Alpine
        guest may own devices only under the control of a confined VMM PD.

  \item \textbf{Exclusive device ownership.}  
        The platform console (UART or virtio-console) shall be owned exclusively
        by a dedicated Console PD. Storage and network devices shall be assigned
        either to native driver PDs or to the Alpine guest; they shall never be
        shared ambiently.

  \item \textbf{MCS scheduling and affinity.}  
        Each PD shall receive an explicit budget, period, and core affinity.
        Priority ordering and resource parameters shall be fixed at build time.

  \item \textbf{Minimal guest images.}  
        Any Linux (or NetBSD Rump) guest shall be constructed from a minimal,
        musl-based root filesystem. Only devices explicitly assigned to the guest
        shall appear in its device tree.

  \item \textbf{Isolated POSIX applications.}  
        The Binary Editor (BE), Terminal Vision (TV), formal POSIX shell (SH),
        and SKYNET swarm coordinator shall each execute as a distinct Protection
        Domain with least-privilege capabilities (see \S\ref{sec:posix-pds}).

  \item \textbf{Zero-dependency userland tools.}  
        BE, TV, and SH shall be implemented in ANSI C99 using only POSIX and
        termios interfaces; external libraries (ncurses, S-Lang, etc.) are
        prohibited.

  \item \textbf{Formal verification preservation.}  
        The seL4 configuration used by Infosec OS.1 shall retain a verified
        status (functional correctness and integrity). The SH PD shall be
        derived from a formal model where claimed.

  \item \textbf{Measurable static image.}  
        The final system image shall support measurement and attestation of its
        static configuration and code.
\end{enumerate}

These requirements are both necessary precursors of the architecture and
directly matched by the current Synrc Hypervision design and its Infosec OS.1
extensions.

\subsection{NIST SP 800-53 Controls}
\label{sec:nist}

Selected control families of NIST SP 800-53 are mapped onto seL4 capabilities
and the static compile-time configuration as follows.

\begin{description}[style=nextline,leftmargin=1.2cm]
  \item[AC --- Access Control]  
        Least privilege and information-flow enforcement are realised by the
        explicit capability layout (AC-3, AC-4, AC-6). Separation of duties is
        obtained by placing distinct functions into distinct PDs.

  \item[SC --- System and Communications Protection]  
        Boundary protection (SC-7) is provided by seL4 Protection Domains and
        stage-2 translation for the Alpine guest. Information-flow control follows
        from the seL4 proofs. All inter-domain communication occurs through
        declared channels only.

  \item[CM --- Configuration Management]  
        The Microkit system description together with compile-time flags
        constitutes the baseline configuration (CM-2, CM-3, CM-6, CM-7). No
        runtime reconfiguration of the capability space is possible.

  \item[AU --- Audit and Accountability]  
        The Monitor PD collects health, performance, and audit-style events via
        fixed shared metrics regions and notifications, supporting AU-2, AU-3,
        AU-6, and AU-12. When MONITOR=0 the domain may be omitted, further
        reducing the TCB.

  \item[SI --- System and Information Integrity]  
        Software integrity (SI-7) is supported by the static, measurable image
        and by the formal verification of the seL4 core.

  \item[CP --- Contingency Planning]  
        Fault isolation is obtained from both BEAM supervision trees and seL4
        spatial/temporal isolation. MCS budgets prevent resource-exhaustion
        interference between domains.
\end{description}

Additional high-assurance properties beyond the NIST catalogue include the
absence of ambient authority, deterministic resource allocation, and the
ability to omit non-essential PDs from production images.

\subsection{POSIX Applications as Protection Domains}
\label{sec:posix-pds}

Elevating BE, TV, SH, and SKYNET into their own Protection Domains constitutes
a deliberate isolation upgrade that applies the same static capability model
already used for Synrc, the VMM, and the device PDs.

\begin{itemize}[leftmargin=*]
  \item \textbf{Spatial isolation.}  
        Each tool receives a private address space. A compromise in the Binary
        Editor (while processing untrusted binaries) cannot reach the memory of
        the shell, the text editor, the BEAM runtime, or the SKYNET control
        plane.

  \item \textbf{Least-privilege capability surfaces.}  
        Capabilities are granted individually:
        \begin{itemize}
          \item BE obtains mediated storage channels and a console channel;
          \item TV obtains only a console channel;
          \item SH obtains a console channel and a tightly restricted filesystem
                view;
          \item SKYNET obtains the network/radio channels required by the swarm
                protocol and a high-priority MCS context.
        \end{itemize}
        No domain receives ambient authority.

  \item \textbf{Fault and temporal isolation.}  
        Crashes, infinite loops, or resource exhaustion remain confined.
        Independent MCS budgets allow SKYNET to run with real-time parameters
        while the interactive tools receive lower-priority contexts.

  \item \textbf{Containment of untrusted input.}  
        Malicious binaries (BE), scripts (SH), free-form text (TV), and external
        swarm messages (SKYNET) are processed inside dedicated domains whose
        outward channels are minimised.

  \item \textbf{Improved auditability.}  
        The Monitor PD can receive distinct notification channels from each
        application PD, producing attributable events rather than mixed
        userland logs.
\end{itemize}

Console multiplexing is performed by the existing Console PD, which continues
to own the physical UART exclusively and exposes virtual character streams.
Storage access is mediated through the Storage PD or a small helper domain.
The NuttX-based real-time layer used by SKYNET may itself reside inside the
SKYNET PD or communicate with it via declared channels.

The net result is a system that more closely approximates a high-assurance
separation kernel: a small number of carefully interfaced application domains
running on a formally verified microkernel with a minimal TCB under the
primary workload.

%----------------------------------------------------------------------
\section{Conclusion}
%----------------------------------------------------------------------

Infosec OS.1 demonstrates that a production-grade BEAM application can operate
with a dramatically reduced trusted computing base while retaining practical
device support and a coherent suite of Infosec and real-time components.
The architecture rests on three pillars: the formally verified seL4 microkernel,
a purpose-built minimal host (Synrc), and capability-mediated confinement of
drivers inside an Alpine guest.

Technical requirements have been stated in accordance with ISO/IEC/IEEE
29148:2018 and are directly realised by the static Microkit construction.
Selected NIST SP 800-53 control families map cleanly onto seL4 capabilities
and compile-time configuration. Elevating the principal POSIX applications
(BE, TV, SH, SKYNET) into dedicated Protection Domains further raises the
assurance level by enforcing spatial isolation, least privilege, fault
containment, and clearer audit trails---without abandoning the static model
or introducing dynamic authority.

The design therefore offers a practical route toward higher-assurance,
certifiable distributed and embedded systems that combine the productivity of
Erlang/OTP with the isolation guarantees of a separation kernel.

\end{document}
